% !TeX root = ../main.tex
% sections/chap3/sec_overview.tex

\subsection{モデルの概要}
\label{sec:overview}
本章ではシミュレーション分析の対象となるニュー・ケインジアン・モデルを構築する。
このモデルは以下の特徴をもつ。
まず経済主体は家計のみであり、家計は生産者を兼ねるヨーマン・ファーマであるとする。
家計は合理的経済人であり、さらに合理的期待形成を仮定する。
国内債券市場については自国と外国それぞれにおいて完備性を仮定する。
国際債券市場については不完備性を仮定する。
本章の構成は以下の通りである。
まず\ref{sec:definitions}節では、モデルの基礎となる定義を記述する。変数の定義や扱い方などが含まれる。
\ref{sec:assumptions}節では、モデルの基礎となる仮定を記述する。
\ref{sec:households}節では、家計の（価格を除く）最適化問題を定式化し1階の条件（FOC）を導出する。
\ref{sec:risk_sharing}節では、国内債券市場を完備としたことの帰結として完全なリスク共有が実現され、
各家計の可処分所得が同一よって消費が同一となることを示す。
\ref{sec:producers}節では、家計が生産者として直面する価格設定問題を取り上げ、
カルボ型の価格硬直性を仮定した下での生産関数と価格分散の動学方程式を導く。
\ref{sec:policies}節では、政府の財政政策および本稿の比較対象となる11の金融政策を定義する。
\ref{sec:resource_constraints}節では、国全体の資源制約式を導出する。
最後に\ref{sec:final_equations}節でこれらの方程式をもとに代表的家計の最終的な方程式体系を得る。